\chapter{Declaración de uso de herramientas de inteligencia artificial}
\label{anexo:ia}

% =====================================================================================
% ⚠️ ANEXO OBLIGATORIO SI SE HAN USADO HERRAMIENTAS DE IA.
%
% «Guía del Trabajo Fin de Estudios», apartado 4.4 «Actividades y entregas», PAGINA 22:
%
%   «...en caso de usar alguna de estas herramientas, el director de TFE debe validar
%    su uso adecuado y ético según las recomendaciones de la universidad. Igualmente,
%    dicho uso debe quedar explícitamente reflejado en el documento incluyendo las
%    referencias específicas a las herramientas de IA utilizadas según las normas de
%    citación oportunas. En el caso de que la universidad detecte un uso inadecuado de
%    herramientas de IA, o una citación inadecuada de las mismas, podrá NO AUTORIZAR el
%    trabajo y declararlo no apto para su defensa. En este caso, podrás perder la
%    convocatoria en curso.»
%
% 🔴 DOS COSAS PENDIENTES, Y LAS DOS SON BLOQUEANTES:
%    1. VALIDACION DEL DIRECTOR. La norma la exige de forma explicita. Hay que
%       planteárselo en sesión y dejar constancia.
%    2. REFERENCIAS APA de cada herramienta en bibliografia.bib, y citarlas desde aquí.
%       Estan añadidas pero SIN VERIFICAR: falta confirmar el nombre exacto del producto,
%       la version y el periodo de uso. No inventarlos.
%
% ⚠️ La declaracion tiene que ser COMPLETA. Omitir una categoria de uso es justo lo que
%    la norma castiga. Si aparece un uso nuevo, se añade aqui.
% =====================================================================================

\section{Alcance de la declaración}
\label{sec:ia_alcance}
% Parrafo de entrada: por que existe este anexo y a que se compromete.
Este anexo detalla, de forma exhaustiva, en qué tareas se han empleado herramientas de
inteligencia artificial generativa durante la realización del trabajo, en cumplimiento de
lo establecido en la guía de la asignatura.

\section{Herramientas empleadas y uso concreto}
\label{sec:ia_herramientas}

\begin{table}[htbp]
    \centering
    \begin{tabular}{@{}p{0.28\textwidth}p{0.62\textwidth}@{}}
        \toprule
        \textbf{Herramienta} & \textbf{Uso} \\
        \midrule
        Claude Code (Anthropic)
            & Desarrollo del código de generación y simulación de circuitos cuánticos
              con Qiskit, y del resto de código ajeno al aprendizaje automático. \\[0.4em]
        Claude Code (Anthropic)
            & Generación de las figuras e ilustraciones de la memoria. \\[0.4em]
        Claude Code (Anthropic)
            & Asistencia en la redacción, la estructuración y la revisión del documento. \\[0.4em]
        Google Gemini
            & Búsqueda bibliográfica en profundidad para la elaboración del estado del
              arte. \\
        \bottomrule
    \end{tabular}

    \caption{Herramientas de inteligencia artificial empleadas y tareas en las que se han
             utilizado.}
    \label{tab:ia}

\end{table}

% 🔴 LAS DOS HERRAMIENTAS SE CITAN AQUI CON \cite{}, y es OBLIGATORIO: la norma pide «las
%    referencias específicas a las herramientas de IA utilizadas según las normas de
%    citación oportunas». Una tabla sola no basta. Las entradas del .bib siguen el formato
%    APA 7 para IA generativa: autor, año, producto y versión, descriptor entre corchetes
%    y URL.

Las dos herramientas quedan referenciadas según la normativa APA que exige la titulación:
\textbf{Claude Code} \cite{anthropic_claude} y \textbf{Gemini Pro}
\cite{google_gemini}.

Las figuras de la memoria son de \textbf{dos clases distintas}, y su pie lo refleja. Las
gráficas que representan medidas del trabajo se han producido con código propio sobre los
datos generados para esta memoria, y su fuente es elaboración propia; los \textbf{diagramas
conceptuales}, que no representan ningún dato, se han elaborado con Claude Code y llevan esa
herramienta como fuente.